% Generated by tools/update_ai_disclosure.py; edit provenance/ai-use.yml instead.
\section*{Declaration of generative AI and AI-assisted technologies}

From 3 August 2026 through the latest covered commit on 19 September 2026, the author used OpenAI Codex (model versions were not consistently recorded), GitHub Copilot (model versions were not consistently recorded), and OpenClaw (model versions were not consistently recorded) within an author-directed, version-controlled manuscript-development workflow. AI assistance was used for proposing and critiquing prose; examining notation and equation references; auditing derivational consistency; identifying candidate literature and bibliographic issues; supporting LaTeX quality assurance; assisting with reproducibility and scientific-software development. The author determined the scientific questions, theoretical structure, assumptions, mathematical arguments, physical interpretations, and final wording. Every accepted change was reviewed by the author; equations and citations were checked against manuscript source and primary literature as applicable, and computational changes were subjected to the repository's applicable validation procedures. AI output was not treated as a scholarly source or credited with authorship. The author takes full responsibility for the accuracy, originality, and integrity of the work.

The repository at \url{https://github.com/johntfoster/finite-strain-biot-poromechanics} is available for inspection of the manuscript's changes and development history.
The repository commit workflow calls for sanitized summaries of the relevant development conversations, including decisions, alternatives, unsuccessful approaches, open questions, and next steps. The commit-message hook requires this structured narrative, and the pre-commit hook regenerates the AI-use disclosure. These records preserve a summary of the development history; the hooks do not collect or archive verbatim chats, and the public repository does not contain a complete prompt history.
